%%%%%%%%%%%%%%%%%%%%%%%%%%%%%%%%%%%%%%%%%%%%%%%%%
%%%%%%%%%%%% chap: architecture %%%%%%%%%%%%%%%%%
%%%%%%%%%%%%%%%%%%%%%%%%%%%%%%%%%%%%%%%%%%%%%%%%%

\chapter{Application Architecture}\label{chapter:architecture}

\section{Architectural Overview}

NewGEO is implemented as a layered monorepo with four main runtime parts:
\begin{itemize}
    \item a shared Python domain package in \texttt{packages/newgeo\_core/};
    \item a FastAPI transport layer in \texttt{api/};
    \item a polling background worker in \texttt{workers/};
    \item a Next.js dashboard in \texttt{web/}.
\end{itemize}

This structure reflects the main design principle of the application: business logic should live
in the shared service layer, not inside HTTP routes, UI components, or worker-specific code.
As a result, the API, tests, scripts, and worker all reuse the same orchestration behavior.

\section{Layered Design}

\subsection{Presentation Layer}

The frontend is implemented as a Next.js application with route-based views for overview,
pages, queries, runs, recommendations, and comparison. The dashboard shell provides a single
navigation and project-selection frame, while the individual pages expose different slices of the
same project dashboard state.

The presentation layer does not contain optimization logic. Its responsibility is to show project
state, benchmark outcomes, recommendation summaries, and review data in a form that is easy
to inspect. This makes it suitable for human-in-the-loop workflows.

\subsection{Transport Layer}

The API layer exposes the application through HTTP endpoints such as:
\begin{itemize}
    \item project creation and dashboard retrieval;
    \item page import and crawl requests;
    \item context-pack and query-cluster creation;
    \item benchmark run creation and artifact retrieval;
    \item recommendation creation, approval, and export;
    \item demo seeding.
\end{itemize}

Each route validates request payloads and delegates to \texttt{NewGeoService}. This keeps the
transport layer thin and prevents HTTP-specific code from becoming the true domain layer.

\subsection{Domain and Service Layer}

The heart of the application is \texttt{packages/newgeo\_core/newgeo\_core/}. This package
contains:
\begin{itemize}
    \item domain models and enums;
    \item storage abstractions;
    \item content normalization utilities;
    \item retrieval helpers for supporting documents;
    \item benchmark connectors and metrics;
    \item recommendation generation and constraint checks;
    \item the high-level orchestration service.
\end{itemize}

The service layer is responsible for the full workflow: page ingestion, context-pack creation,
query-cluster creation, benchmark execution, recommendation generation, approval, export,
dashboard aggregation, and queued-job processing.

\subsection{Background Execution Layer}

The worker process repeatedly polls the store for queued runs and queued recommendations. If
it finds any, it executes them using the same service-layer methods used by synchronous API
calls. Although this worker is intentionally simple, it proves the architectural boundary between
request submission and deferred execution.

\section{Domain Model}

The application models the optimization workflow through explicit entities rather than implicit
state transitions.

\subsection{Project and Page Model}

A \texttt{Project} is the top-level monitored entity. It groups all pages, context packs, query
clusters, runs, and recommendations related to a content property. A \texttt{Page} represents a
tracked document, while a \texttt{ContentSnapshot} stores the actual normalized body and its
derived embedding. This separation is useful because a page identity can remain stable while its
content evolves over time.

\subsection{Context Model}

A \texttt{ContextPack} stores the external information used to guide rewriting. It includes a
brief, voice rules, supporting documents, and typed constraints. Constraints may represent
locked facts, forbidden claims, required terms, or voice guidance. This model is the central
mechanism by which external data influences the recommendation pipeline in a controlled way.

\subsection{Query and Benchmark Model}

A \texttt{QueryCluster} groups related user queries under one topical unit. This is important
because visibility is query-dependent. A page may perform well for one demand pattern and
poorly for another.

Each benchmark execution is represented by a \texttt{BenchmarkRun}. The run stores its engine
name, run type, prompt version, connector configuration, candidate pages, observations, score
snapshots, and persisted artifacts. Individual scored outcomes are stored as
\texttt{ScoreSnapshot} records so trends can be reconstructed later.

\subsection{Recommendation and Approval Model}

A generated suggestion is represented by a \texttt{Recommendation}. The actual draft output is
stored in a \texttt{RecommendationBundle}, which includes rewritten markdown, a markdown
diff, rationale, confidence, supporting context used, constraint checks, and a preview of
baseline versus projected benchmark metrics.

Human review is represented explicitly through the \texttt{Approval} model. This keeps the
publication workflow auditable and avoids silently treating generated text as approved text.

\section{Data Flow}

\subsection{Benchmark Flow}

The benchmark execution flow proceeds as follows:
\begin{enumerate}
    \item the user creates a run through the API;
    \item the service resolves the query cluster and candidate pages;
    \item a connector is selected based on run type and connector configuration;
    \item the connector scores candidate pages and returns observations and artifacts;
    \item the service stores score snapshots and artifact records;
    \item the dashboard later aggregates these results for visualization.
\end{enumerate}

This flow is designed so that raw prompt/response material and aggregate metrics are both
preserved. That separation is important for future debugging and connector replacement.

\subsection{Recommendation Flow}

The recommendation flow is similar, but includes more safety logic:
\begin{enumerate}
    \item the user requests a recommendation for one page, one context pack, and one query cluster;
    \item the service loads the current page snapshot and relevant candidate pages;
    \item supporting documents are ranked against the query and brief;
    \item a conservative rewrite draft is assembled;
    \item constraint checks and source-claim preservation checks are applied;
    \item the original and projected versions are benchmarked;
    \item the recommendation bundle is stored for review.
\end{enumerate}

Approval and export then form a separate downstream stage. This prevents generation from
being confused with publication.

\section{Storage Architecture}

The storage layer is defined through a \texttt{StorageBackend} protocol with methods for saving,
loading, listing, and deleting models. Two concrete implementations are available:
\begin{itemize}
    \item \texttt{JsonFileBackend}, used as the default local-development store;
    \item \texttt{SqliteStore}, which persists the same logical collections in a relational table.
\end{itemize}

The service itself uses a \texttt{JsonStore} compatibility facade that resolves to the selected
backend. This design choice keeps the service logic stable while allowing local JSON storage,
SQLite migration, and future replacement with a PostgreSQL-based adapter.

From an architectural standpoint, this is an important compromise. JSON storage keeps the
prototype frictionless and easy to run. SQLite demonstrates that the service logic can already
survive a backend swap. The storage protocol therefore acts as a real extension boundary rather
than a hypothetical one.

\section{Frontend State and Project Awareness}

The frontend resolves a selected project through query parameters and uses that project to load
the dashboard state. If a live API base URL is not configured, it falls back to sample data. This
choice allows the interface to remain explorable even before all mutation flows are fully wired.

Architecturally, the UI is therefore ahead of full integration but not disconnected from the real
backend shape. The route structure, project-aware URLs, dashboard aggregation format, and
comparison views all follow the live domain model. This makes the current frontend a practical
product shell rather than a disposable mockup.

\section{Architectural Rationale}

The application architecture is intentionally more modular than a small prototype strictly
requires. This is justified by the thesis topic itself. Because the system studies how external
data shapes model outputs, it must make those transformation points visible:
\begin{itemize}
    \item content normalization is isolated;
    \item retrieval of supporting context is isolated;
    \item benchmark connectors are isolated;
    \item recommendation generation is isolated;
    \item approval is isolated from generation;
    \item persistence is isolated behind a backend protocol.
\end{itemize}

This separation makes the influence surface inspectable. It also makes the prototype suitable as
the basis for future connectors, richer storage, stronger semantic checks, and real deployment
integrations.
